\scanpage{169}
\mspage{118}
% archive sheet number: 169
La proposition 12 nous donne donc la formule

\[
\pi_0({}_y\backslash X)=\underset{I}{\unsure{\underline{\lim}|\varinjlim}}\,\pi_0({}_y\backslash X_i),
\]
laquelle formule ne peut s'empêcher d'être fonctorielle en $y$, \unsure{tout} comme celle de prop.~12 \deleted{\unsure{le}} l'est dans $\text{\uline{Prop}}_W\,I$. Or \deleted{\unsure{faisons}}\inserted{pour} $y$ variable dans $Y$, le \deleted{foncteur} préfaisceau sur $Y$

\[
y\longmapsto\pi_0({}_y\backslash X_i)
\]
n'est autre que $p_{i\,!}(e_{X_i})$, où

\[
p_i:X_i\longrightarrow Y
\]
est le composé $X_i\to X\xrightarrow{p}Y$ (induit par $p$).

On a ainsi prouvé

\uline{Corollaire} 1 (de prop.~12) Soit $f:X\to I$ $W$-propre,
% slanted
\margin{dire : $W_0$-propre \ldots} et $p:X\to Y$ un foncteur quelconque. Soit, pour tout $i\in\mathrm{Ob}\,I$, $p_i:X_i\to Y$ le foncteur induit sur la fibre \inserted{$X_i$} de $X$ en $i$. On a alors une loi \uline{fonctorielle} naturelle

\[
i\longmapsto p_{i\,!}(e_{X_i}):\quad I\longrightarrow Y^\wedge,
\]
et on a la formule

\[
p_!(e_X)\xleftarrow{\sim}\underset{I}{\varinjlim}\,p_{i\,!}(e_{X_i}),\tag*{($\alpha$)}
\]
où la flèche est induite par les morphismes canoniques

\[
p_{i\,!}(e_{X_i})\to p_!(e_X)
\]
